\section{附录：符号与术语表}\label{sec:appendix}

\begin{table}[htbp]
    \centering
    \caption{本文使用的主要符号}
    \label{tab:notation}
    \begin{tabular}{ll}
        \toprule
        符号 & 含义 \\
        \midrule
        $E$ & Besicovitch 集（挂谷集） \\
        $E_\delta$ & $E$ 的 $\delta$-邻域 \\
        $\delta$-管 & 长 $1$、粗 $\delta$ 的圆柱形邻域 \\
        $S^{n-1}$ & $\mathbb{R}^n$ 中的单位球面（方向空间） \\
        $\dim_{\mathrm{H}}$ & Hausdorff 维数 \\
        $\dim_{\mathrm{M}}$ & Minkowski（盒计数）维数 \\
        $M_\delta$ & Kakeya 极大算子 \\
        $\mathbb{T}_j$ & 第 $j$ 组 $\delta$-管族（横向） \\
        $\chi_{T}$ & 管子 $T$ 的示性函数 \\
        $|K|$ & 集合的测度（有限域下为基数） \\
        $\mathbb{F}_q$ & $q$ 元有限域 \\
        $\gtrsim$ & 常数倍不等式（$A \gtrsim B$ 即 $A \ge cB$） \\
        $f \mid_T$ & $f$ 在 $T$ 上的平均 \\
        \bottomrule
    \end{tabular}
\end{table}

\section*{术语对照}
\begin{itemize}
    \item 转针问题（Kakeya needle problem）：求使单位线段连续旋转 $180^\circ$ 的最小面积区域；
    \item Besicovitch 集 / 挂谷集（Besicovitch set / Kakeya set）：包含所有方向单位线段的（测度为零的）紧集；
    \item Perron 树（Perron tree）：通过平移细分三角形大幅交叠而得到的低面积构造；
    \item 灌木丛论证（bush argument）：以共点管族的 $L^2$ 正交性计数的技术；
    \item 毛刷论证（hairbrush argument）：Wolff 的按角度收集管族、分情形计数的技术；
    \item 多线性 Kakeya 估计（multilinear Kakeya estimate）：横向管族并集的体积下界；
    \item 多项式分划（polynomial partitioning）：用多项式零点集剖分空间以分离代数与非代数行为的方法；
    \item 多项式方法（polynomial method）：以低次多项式零点整除结构导出组合下界的方法；
    \item 限制性估计（restriction estimates）：Fourier 变换在曲面上限制的 $L^p$ 界；
    \item 局部光滑化（local smoothing）：色散方程解的时空正则性增益猜想；
    \item Nikodym 集（Nikodym set）：过每点存在一条仅与集合交于该点的直线的集合。
\end{itemize}
